\subsection*{EB valence is shaped by context, sex, and circuit composition}
To test whether the valence effects observed in the FLIC generalized beyond a feeding context, we assayed place preference in a rectangular multi-lane arena in which one half of each lane was illuminated and the other remained dark. Unlike the closed-loop FLIC, this open-loop design delivered light on a pre-programmed schedule \hl{[specify cycle pattern]} regardless of fly position, dissociating valence from sucrose contact or response-contingent stimulation. Flies were observed during a 10-min dark acclimation, a 60-min treatment, and a 10-min dark recovery; persistence of side bias into recovery was used to distinguish stable state effects from acute responses to light. Both sexes were tested, and each manipulation was evaluated at two stimulation frequencies (48 Hz and 200 Hz) as a robustness check.

The positive valence of R2/R4 co-activation generalized cleanly. Both sexes significantly preferred the illuminated side during treatment, and it was observed at both 48 Hz (\figref{fig:arenafig}D) and 200 Hz (\figref{fig:arenasupp}D), indicating that the appetitive state reflects a stable property of the R2/R4-activated circuit rather than an artifact of a particular assay or stimulation regime.

R2 activation also generalized, but with important qualifications. Females exhibited clear avoidance of the illuminated side (\figref{fig:arenafig}B, left), which was reproduced at higher frequency (\figref{fig:arenasupp}A, left) and was consistent with the aversive signature observed in the FLIC. The arena phenotype was, however, more persistent: FLIC females disengaged rapidly from the sucrose well but continued to re-approach at normal frequency, whereas arena females sustained avoidance throughout treatment \hl{I'm not sure I understand this from your original writing}. This contrast suggests that the behavioral expression of R2-dependent aversion, though not the underlying valence, is shaped by assay context, likely by the presence of food, the structure of the environment, or the temporal pattern of stimulation. In males, R2 activation produced the opposite preference: males preferred the illuminated side during both treatment and recovery at 48 Hz (\figref{fig:arenafig}B, right), with a weaker effect at 200 Hz (\figref{fig:arenasupp}A, right). Thus, the arena revealed a sex-dependent valence for R2 activation.

By contrast, the FLIC findings for R4d manipulation did not generalize robustly. R4d>GtACR1 and control flies showed indistinguishable side biases (females preferring illumination, males avoiding it; \figref{fig:arenasupp}B), so although the open-loop design removed the light-startle confound, it did not permit an unambiguous valence assignment. R4d activation, which had been aversive in the FLIC, produced variable behavior in the arena: females showed little avoidance at 48 Hz but avoided illumination at 200 Hz (\figref{fig:arenafig}C, left; \figref{fig:arenasupp}C, left), while males avoided illumination at 48 Hz but showed no preference at 200 Hz (\figref{fig:arenafig}C, right; \figref{fig:arenasupp}C, right). The absence of a stable signature across conditions suggests that R4d neurons act in a more context-sensitive or modulatory fashion than R2 or R2/R4.

In summary, the behavioral consequences of EB manipulation are nonlinear and depend on circuit composition, sex, and assay context. R2/R4 co-activation produced a robust positive valence that generalized across feeding-dependent and feeding-independent paradigms and persisted after chronic stimulation, indicating a stable property of circuit state rather than an acute stimulation artifact; in the FLIC, this state was expressed as more frequent returns to the sucrose well without prolonged visit duration, consistent with enhanced approach or incentive salience rather than increased consummatory behavior \textit{per se}. R2 activation was consistently aversive in females across both assays but produced the opposite preference in males, revealing a sex-dependent valence. R4d manipulations, by contrast, did not yield a valence signature stable enough to generalize across assays. Together, these findings argue that EB ring neurons do not make fixed, additive contributions to motivational state, but instead generate qualitatively distinct and differentially stable internal states through combinatorial recruitment, raising the question of how such state differences relate to the common downstream outcome of extended lifespan.

